\section{Introduction}\label{sec:intro}

Large language models have demonstrated remarkable mathematical reasoning
capabilities when augmented with chain-of-thought (CoT) prompting, in which a
model produces a sequence of intermediate reasoning steps before arriving at a
final answer~\cite{wei2022chain}.  A natural application is \emph{proof
verification}: given a purported proof of a mathematical statement, a model
checks correctness step by step.  Recent work shows that verification accuracy
can be substantially improved by generating multiple CoT traces and aggregating
their judgments via majority voting---a strategy known as
\emph{self-consistency}~\cite{wang2022self}---and that the returns from
additional test-time compute follow difficulty-dependent saturation
curves~\cite{snell2024scaling, shi2025heimdall}.

A central question, however, remains open: \emph{what structural properties of
a proof determine the difficulty of verifying it?}  Snell et
al.~\cite{snell2024scaling} characterize difficulty using the model-dependent
pass@1 rate, while Shi and Jin~\cite{shi2025heimdall} demonstrate that
verification difficulty does not simply track problem-solving difficulty.
Neither framework identifies proof-intrinsic features that predict verification
convergence a priori.

Proof complexity theory offers a rich collection of such features.  The
\emph{length} (number of inference steps), \emph{width} (maximum clause size),
\emph{depth} (height of the proof DAG), and \emph{space} (maximum
simultaneously stored clauses) of resolution proofs are well studied and
exhibit deep interrelationships~\cite{bensasson2001short, haken1985intractability}.
Jarvisa\"alo et al.~\cite{jarvisalo2012relating} showed that space complexity
correlates with practical SAT solver running time more strongly than length or
width.  Yet no prior work connects these measures to LLM verification behavior.

In this paper we bridge proof complexity theory and CoT verification
convergence.  We model CoT verification as a stochastic saturation process
whose parameters are governed by proof structure, and we prove rigorous bounds
relating proof complexity to verification accuracy and required compute.  Our
contributions are as follows.

\begin{itemize}[leftmargin=*,itemsep=2pt,topsep=4pt]
\item We formalize CoT verification convergence as an exponential saturation
  process and prove that the convergence step scales as
  $O(\log(1/\varepsilon)/\lambda)$, where $\lambda$ depends inversely on a
  composite proof complexity measure (Theorem~\ref{thm:convergence}).
\item We establish an error accumulation bound showing that asymptotic
  verification accuracy decays exponentially in proof length under a
  step-independence assumption (Theorem~\ref{thm:error-accumulation}).
\item We prove that self-consistency amplification via majority voting requires
  $\Theta(1/(p - \tfrac{1}{2})^2)$ independent chains, connecting proof
  complexity to test-time compute scaling
  (Theorem~\ref{thm:self-consistency}).
\item We validate these predictions computationally on 129~synthetic proof
  instances across four formula families, obtaining strong rank correlations
  ($\rho = -0.99$) between the composite complexity measure and convergence
  rate.  We identify a phase transition in verification difficulty at a critical
  complexity threshold (Proposition~\ref{prop:phase-transition}).
\end{itemize}

The paper is organized as follows.  Section~\ref{sec:prelim} introduces the
necessary background from proof complexity theory and CoT reasoning.
Section~\ref{sec:main} presents our main results with complete proofs and
computational verification.  Section~\ref{sec:discussion} discusses connections
to prior work, implications, and open questions.  Section~\ref{sec:conclusion}
concludes.
